%%%%%%%%%%%%%%%%%%%%%%%%%%%%%%%%%%%%%%%%%
% Memo
% LaTeX Template
% Version 1.0 (30/12/13)
%
% This template has been downloaded from:
% http://www.LaTeXTemplates.com
%
% Original author:
% Rob Oakes (http://www.oak-tree.us) with modifications by:
% Vel (vel@latextemplates.com)
%
% License:
% CC BY-NC-SA 3.0 (http://creativecommons.org/licenses/by-nc-sa/3.0/)
%
%%%%%%%%%%%%%%%%%%%%%%%%%%%%%%%%%%%%%%%%%

\documentclass[letterpaper,11pt]{texMemo} % Set the paper size (letterpaper, a4paper, etc) and font size (10pt, 11pt or 12pt)

\usepackage{fancyhdr}
\usepackage{fancybox}
\usepackage{longtable}
\usepackage{amsmath}
\usepackage{xcolor}
%----------------------------------------------------------------------------------------
%	MEMO INFORMATION
%----------------------------------------------------------------------------------------



\memoto{Luis Andr\'es Valido Fajardo. luis.valido@umcc.cu (53694742)} % Recipient(s)

\memofrom{Josval Díaz Blanco} % Sender(s)

\memosubject{Guía de Aprendizaje para Concursantes ICPC y IOI: Búsqueda Binaria } % Memo subject

\memodate{\today} % Date, set to \today for automatically printing todays date

\logo{\includegraphics[scale=0.5]{img/icpc}} % Institution logo at the top right of the memo, comment out this line for no logo

%----------------------------------------------------------------------------------------

%\titleguide{Detección de ciclos en un grafo}%cycle_graph
%\titleguide{Camino y ciclo de Hamilton}%hamilton_tour
%\titleguide{Encontrarse en el medio ( \emph{Meet in the Middle} )}%meet_in_middle
%\titleguide{Números Catalanes}%number_catalan
%\titleguide{Introducción al ajedrez}%introduction_chess 
%\titleguide{Operaciones con matrices}%matrix_operation 
%\titleguide{Coeficientes binomiales}%binomial_coefficients
%\titleguide{Funciones de distancias}%_distance_functions
%\titleguide{Tabla dispersa (\emph{Sparse Table})}%sparse_table 
%\titleguide{Algoritmo-Z (\emph{Z-Algorithm})}%z_algorithm 
%\titleguide{Función de prefijo. Algoritmo de Knuth-Morris-Pratt (\emph{KMP})}%kmp
%\titleguide{Descomposición raíz cuadrada (\emph{Sqrt Decomposition})}%sqrt_decomposition
%\titleguide{Algoritmo de Mo (\emph{Mo's algorithm})}%mo_algorithm 
%\titleguide{Probabilidades}%probability
%\titleguide{Funciones de String en C++}%f_string_cpp
%\titleguide{Búsqueda ternaria}%_ternary_search
%\titleguide{El método de Newton para encontrar raíces}%_roots_newton
%\titleguide{Algoritmo Extendido de Euclides}%extended_euclidean_algorithm
%\titleguide{Ecuación diofántica lineal}%_linear_diophantine_equation
%\titleguide{Sucesión de Fibonacci}%_fibonacci
%\titleguide{Estructuras de datos basadas en políticas (árbol de estadísticas de pedidos)}%_policy_based_data_strctures
%\titleguide{Exponenciación binaria}%_pow_binary
%\titleguide{Algoritmo de la criba Eratóstenes}%sieve_eratosthenes
%\titleguide{Algoritmo de la criba linear}%_linear_sieve
%\titleguide{Pruebas de primalidad}%_primality_test
%\titleguide{Factorización de enteros}%_integer_factorization
%\titleguide{La función totiente de Euler ($\phi(n)$)}%_euler_totient_function
%\titleguide{Arreglos de sufijos (\emph{Suffix Array})}%suffix_array FALTA
%\titleguide{Permutaciones}%permutation FALTA
%\titleguide{Sucesiones númericas}%successions FALTA
%\titleguide{Arbol de raíz cuadrada (\emph{SQRT Tree})}%sqrt_tree FALTA
%\titleguide{Cadenas de carácteres, terminología}%_string_terminology FALTA
\titleguide{Cantidad y suma de los divisores}%_numbers_and_sums_divisors
\begin{document}

%\AddToShipoutPicture{\BackgroundPic}
\maketitle % Print the memo header information
%\tableofcontents
\pagestyle{plain}
\pagebreak

\pagestyle{fancy}
\fancyhead[LO,CE]{ }
\fancyhead[RO,CE]{\includegraphics[scale=0.1]{img/icpc}}
\fancyfoot[LO,CE]{\textbf{Autor:} Luis Andrés Valido Fajardo \\ \textbf{Email:} luis.valido1989@gmail.com \\ \textbf{Teléfono:} 53694742}
\fancyfoot[RO,CE]{\emph{Existen 10 tipos de personas Las que \\saben binario y LAS QUE NO}}
\fancypagestyle{plain}{\pagestyle{fancy}}



%\lhead{ }
%\rhead{  }

%\fancyfoot[L]{}
%\fancyfoot[R]{\textbf{Autor:} Luis Andrés Valido Fajardo \\ \textbf{Email:} luis.valido@umcc.cu}
%----------------------------------------------------------------------------------------
%	MEMO CONTENT
%----------------------------------------------------------------------------------------


\section{Introducción}
En algunos problemas o ejercicios vinculados a la teoría de números el elemento clave de la solución radica en hallar o bien la cantidad de divisores de un número o la suma de estos. A como hallar ambas cuestiones de una manera eficiente lo veremos a continuación.
\section{Conocimientos previos}
\subsection{Número primo}
En matemáticas, un número primo es un número natural mayor que 1 que tiene únicamente
dos divisores positivos distintos: él mismo y el 1. Por el contrario, los números compuestos son
los números naturales que tienen algún divisor natural aparte de sí mismos y del 1, y, por lo tanto, pueden factorizarse. El número 1, por convenio, no se considera ni primo ni compuesto

\subsection{Factorización en números primos}
Los factores primos de un número entero son los números primos divisores exactos de ese número entero. El proceso de búsqueda de esos divisores se denomina factorización de enteros, o factorización en números primos.
\section{Desarrollo}
\subsection{Números de divisores}

Debe ser obvio que la factorización principal de un divisor $d$ tiene que ser un subconjunto de la factorización principal de $n$ , por ejemplo. $6 = 2 \cdot 3$ es un divisor de $60 = 2^2  \cdot 3  \cdot 5$ . Así que sólo se está consolviendo todos los subconjuntos diferentes de la factorización principal de $n$.

Por lo general, el número de subconjuntos es de $2^x$ por un conjunto con $x$ elementos. Sin embargo, esto ya no es cierto, si hay elementos repetidos en el conjunto. En nuestro caso, algunos factores principales pueden aparecer varias veces en la factorización principal de $n$.

Si un factor primo $p$ aparece $e$ veces en la factorización principal de $n$, entonces podemos usar el factor $p$ hasta la $e$ veces en el subconjunto. Lo que significa que tenemos $e+1$ opciones.

Por lo tanto, si la factorización principal de $n$ es $p_1^{e_1} \cdot p_2^{e_2} \cdots p_k^{e_k}$, donde $p_i$ son números primos diferentes entonces, el número de divisores es:  

$$d(n) = (e_1 + 1) \cdot (e_2 + 1) \cdots (e_k + 1)$$

Una forma de pensarlo es lo siguiente:

\begin{itemize}
	\item Si hay un único divisor primo $n = p_1^{e_1}$, entonces obviamente tendrá $e_1 + 1$ divisores ($1, p_1, p_1^2, \dots, p_1^{e_1}$).
	\item Si hay dos divisores primos diferentes  $n = p_1^{e_1} \cdot p_2^{e_2}$, entonces se puede organizar todos los divisores de forma tabular.
	
	$$\begin{array}{c|ccccc} & 1 & p_2 & p_2^2 & \dots & p_2^{e_2} \\\\\hline 1 & 1 & p_2 & p_2^2 & \dots & p_2^{e_2} \\\\ p_1 & p_1 & p_1 \cdot p_2 & p_1 \cdot p_2^2 & \dots & p_1 \cdot p_2^{e_2} \\\\ p_1^2 & p_1^2 & p_1^2 \cdot p_2 & p_1^2 \cdot p_2^2 & \dots & p_1^2 \cdot p_2^{e_2} \\\\ \vdots & \vdots & \vdots & \vdots & \ddots & \vdots \\\\ p_1^{e_1} & p_1^{e_1} & p_1^{e_1} \cdot p_2 & p_1^{e_1} \cdot p_2^2 & \dots & p_1^{e_1} \cdot p_2^{e_2} \\\\ \end{array}$$
	
	lo que nos indica que el número de divisores en este caso es $(e_1 + 1) \cdot (e_2 + 1)$.
	
	\item Se puede argumentar de manera similar si hay más de dos factores primos distintos.
\end{itemize}

\subsection{Suma de divisores}

Podemos realizar el razonamiento de la sección anterior.

\begin{itemize}
	\item Si hay un único divisor primo $n = p_1^{e_1}$, entonces las suma es:
	
	$$1 + p_1 + p_1^2 + \dots + p_1^{e_1} = \frac{p_1^{e_1 + 1} - 1}{p_1 - 1}$$
	
	\item Si hay dos divisores primos diferentes  $n = p_1^{e_1} \cdot p_2^{e_2}$, podemos hacer la misma tabla que antes. La diferencia única es que ahora se puede calcular la suma en el lugar de contar los elementos. Es fácil ver que la suma de cada combinación puede expresarse como:
	
	$$\left(1 + p_1 + p_1^2 + \dots + p_1^{e_1}\right) \cdot \left(1 + p_2 + p_2^2 + \dots + p_2^{e_2}\right)$$
	$$ = \frac{p_1^{e_1 + 1} - 1}{p_1 - 1} \cdot \frac{p_2^{e_2 + 1} - 1}{p_2 - 1}$$
	
	\item En general, para $n = p_1^{e_1} \cdot p_2^{e_2} \cdots p_k^{e_k}$ tenemos la fórmula:
	
	$$\sigma(n) = \frac{p_1^{e_1 + 1} - 1}{p_1 - 1} \cdot \frac{p_2^{e_2 + 1} - 1}{p_2 - 1} \cdots \frac{p_k^{e_k + 1} - 1}{p_k - 1}$$
	
\end{itemize}

\subsection{Funciones multiplicativas}

Una función multiplicativa es una función $f (x)$ que se consina

$$f(a \cdot b) = f(a) \cdot f(b)$$

si $a$ y $b$ son coprimos.

Tanto $d(n)$ y $\sigma (n)$ son funciones multiplicativas.

Las funciones multiplicativas tienen una gran variedad de propiedades interesantes, que pueden ser muy útiles en los problemas de teoría de números. Por ejemplo, la convolución Dirichlet de dos funciones multiplicativas también es multiplicativa.
\section{Implementación}
\subsection{C++}

\subsubsection{Números de divisores}
\begin{lstlisting}[language=C++]
long long number_divisors(long long num) {
   long long total = 1;
   for (long long i = 2; (long long)i * i <= num; i++) {
      if (num % i == 0) {
         long long e = 0;
         do {
            e++;
            num /= i;
         } while (num % i == 0);
         total *= e + 1;
      }
   }
   if (num > 1) total *= 2;
   return total;
}
\end{lstlisting}


\subsubsection{Suma de divisores}
\begin{lstlisting}[language=C++]
long long sum_divisors(long long num) {
   long long total = 1;
   for (long long i = 2; (long long)i * i <= num; i++) {
      if (num % i == 0) {
         long long e = 0;
         do {
            e++;
            num /= i;
         } while (num % i == 0);
         long long sum = 0, pow = 1;
         do {
            sum += pow;
            pow *= i;
         } while (e-- > 0);
         total *= sum;
      }
   }
   if (num > 1) total *= (1 + num);
   return total;
}
\end{lstlisting}


\subsection{Java}

\subsubsection{Números de divisores}
\begin{lstlisting}[language=Java]
public static long number_divisors(long num) {
  long total = 1;
  for (long i = 2; i * i <= num; i++) {
     if (num % i == 0) {
        long e = 0;
        do {
           e++;
           num /= i;
        } while (num % i == 0);
        total *= e + 1;
     }
  }
  if (num > 1) total *= 2;
  return total;
}
\end{lstlisting}


\subsubsection{Suma de divisores}
\begin{lstlisting}[language=Java]
public static long sum_divisors(long num) {
   long total = 1;
   for (long i = 2; i * i <= num; i++) {
      if (num % i == 0) {
         long e = 0;
         do {
            e++;
            num /= i;
         } while (num % i == 0);
           long sum = 0, pow = 1;
         do {
            sum += pow;
            pow *= i;
         } while (e-- > 0);
         total *= sum;
      }
   }
   if (num > 1) total *= (1 + num);
   return total;
}
\end{lstlisting}

\section{Aplicaciones}
Entre las aplicaciones que tiene el poder determinar la cantidad y suma de los divisores de un número $n$ están:

\begin{enumerate}
	\item \textbf{Criptografía:} En criptografía, el cálculo de la cantidad o suma de divisores  de un número puede ser utilizado en la generación de claves criptográficas seguras.
	
	\item \textbf{Teoría de números:} El cálculo de la cantidad de divisores de un número es fundamental en la teoría de números para estudiar propiedades de los números enteros, como los números primos y compuestos. Calcular la suma de los divisores de un número es útil en la teoría de números para estudiar propiedades de los números enteros, como los números perfectos, abundantes y deficientes.
	
	\item \textbf{Optimización de algoritmos:} En informática, el cálculo de la cantidad o suma de divisores de un número puede ser utilizado para optimizar algoritmos en diferentes áreas, como la programación dinámica y la optimización combinatoria.
	
	\item \textbf{Matemáticas recreativas:} El cálculo de la cantidad o suma de divisores de un número también puede ser utilizado en problemas matemáticos recreativos y desafíos matemáticos para ejercitar el razonamiento lógico y matemático. 
	
	\item  \textbf{Ingeniería eléctrica:} En ingeniería eléctrica, el cálculo de la cantidad o suma de divisores de un número puede ser utilizado en el diseño y análisis de circuitos eléctricos para determinar la eficiencia y la estabilidad del sistema.
	
\end{enumerate}

\section{Complejidad}
La complejidad temporal de ambas implementaciones es O($\sqrt{N}\times e$) siendo $N$ el valor al cual le desea determinar sus divisores y $e$ el mayor exponente de los factores primos presentes en la factorización de $N$.
\section{Ejercicios}
A continuación una lista de ejercicios donde se pueden aplicar los contenidos abordados:

\begin{itemize}
	\item \href{https://www.spoj.com/problems/COMDIV/}{SPOJ - COMDIV}
	\item \href{https://www.spoj.com/problems/DIVSUM/}{SPOJ - DIVSUM}
	\item \href{https://www.spoj.com/problems/DIVSUM2/}{SPOJ - DIVSUM2}
\end{itemize}
%\section{Anexos}
%\input{tex_files/anexos_successions.tex}


\end{document}